\section{Research Problem and Scientific Basis}

\subsection{Why IoT traffic can be modelled as counts}

Many IoT devices perform a narrow set of tasks and repeatedly communicate with
the same local and cloud services. Their traffic is not constant, but its protocol
mixture often has a recognisable structure. For one window the system stores

\[
\mathbf{x}=(x_{\mathrm{TCP}},x_{\mathrm{UDP}},x_{\mathrm{SSDP}},x_{\mathrm{ARP}}),
\qquad N=\sum_k x_k.
\]

The categories are mutually exclusive: SSDP is separated from the remaining UDP
traffic. This representation is small enough for an edge-oriented detector and
retains protocol-level evidence that can be explained to an operator.

A fixed multinomial distribution is often too rigid because ordinary device
behaviour changes from window to window. The Dirichlet--multinomial distribution
adds this overdispersion through a positive parameter vector
$\boldsymbol{\alpha}$. Its expected traffic mixture is
$\alpha_k/\alpha_0$, where $\alpha_0=\sum_k\alpha_k$. The project also reports
$\psi=1/\alpha_0$ as a compact dispersion indicator.

The complete log-probability of a window is

\[
\log P(\mathbf{x}\mid\boldsymbol{\alpha})
=\log\frac{N!}{\prod_k x_k!}
+\log\frac{\Gamma(\alpha_0)}{\Gamma(\alpha_0+N)}
+\sum_k\log\frac{\Gamma(\alpha_k+x_k)}{\Gamma(\alpha_k)}.
\]

A low score means that the observed mixture is unlikely under the normal model.
The decision threshold $\tau$ is learned from a separate benign calibration
period:

\[
\text{anomaly if }\log P(\mathbf{x}\mid\boldsymbol{\alpha})<\tau.
\]

The versioned experiments use a raw score and a one-percent lower calibration
quantile. Per-packet normalisation remains available as a separate experiment;
it is not silently mixed with the reported results.

\subsection{Dispersion and parameter estimation}

The model is conditional on the observed total $N$. Writing
$p_k=\alpha_k/\alpha_0$, its conditional variance is
\[
\operatorname{Var}(X_k\mid N)
=N p_k(1-p_k)\frac{N+\alpha_0}{1+\alpha_0}.
\]
For the same expected mixture, smaller concentration permits greater
window-to-window variation. The reported $\psi=1/\alpha_0$ is an inverse
concentration; the common correlation parameter is
$1/(1+\alpha_0)=\psi/(1+\psi)$. It should not be confused with an unrelated
software threshold or probability of attack.

The fitting objective omits only the count-dependent multinomial coefficient:
\[
K(\mathbf{x};\boldsymbol{\alpha})=
\log\Gamma(\alpha_0)-\log\Gamma(N+\alpha_0)
+\sum_k\{\log\Gamma(x_k+\alpha_k)-\log\Gamma(\alpha_k)\}.
\]
For independent rows the objective is $\sum_iK(\mathbf{x}_i;\boldsymbol{\alpha})$.
The complete window score adds
$\log\Gamma(N+1)-\sum_k\log\Gamma(x_k+1)$.
The optional normalised convention divides the complete score by
$\max(N,1)$. Calibration and scoring must select the same convention.
Training kernel values and anomaly thresholds therefore belong to different
scales and cannot be compared directly.

The initial alpha vector is the pooled category proportion multiplied by the
configured initial concentration, usually ten. Minka's fixed-point update is
\[
\alpha_k^{\mathrm{new}}
=\alpha_k\,
\frac{\sum_i[\operatorname{digamma}(x_{ik}+\alpha_k)-\operatorname{digamma}(\alpha_k)]}
{\sum_i[\operatorname{digamma}(N_i+\alpha_0)-\operatorname{digamma}(\alpha_0)]}.
\]
The implementation uses SciPy's digamma for this update regardless of the
likelihood backend. It stops when the absolute likelihood change is below the
configured tolerance, or returns non-convergence at its iteration limit.
Changing the likelihood backend is not replacing this estimator.

\subsection{The role of the LM method}

The Languasco--Migliardi (LM) method evaluates close log-Gamma differences
accurately inside the Dirichlet--multinomial likelihood \cite{lmcore,tpami}.
Parameter estimation uses Minka's
fixed-point equations \cite{minka}; LM or SciPy supplies the likelihood used for
convergence checks and scoring.

This distinction prevents an important overclaim. A numerical backend can change
runtime and numerical behaviour near the multinomial limit, but it does not make
an unchanged statistical detector inherently more accurate. The system keeps a
SciPy implementation as an independent reference and selects six-digit native LM
in the versioned research configurations.

\subsection{Relation to the supplied IoT work}

The reference study models D-Link Camera 5 traffic as daily ten-minute vectors
with the same four categories and compares likelihood methods in modelling and
forecasting experiments \cite{iotpaper}. \projectname{} retains that statistical
foundation but changes the application protocol. It fits one device baseline,
uses a separate threshold-calibration period, emits anomaly decisions, and tests
those decisions on labelled traffic from additional IoT devices.

The D-Link captures remain useful as a matched numerical workload. They do not
contain attack labels in this repository, so they cannot support measured
precision, recall, or F1. This is why the main anomaly results use UNSW Chromecast
and Samsung camera traces instead.

\subsection{What an event explains}

The expected category profile is $\boldsymbol{p}$ and the residual for an
observed count is $x_k/N-p_k$ when $N>0$. For silence the observed proportions
are represented by zeros, so the residual is $-p_k$. Severity is
\[
\frac{\max(0,\tau-s(\mathbf{x}))}{\mathrm{IQR}_{\mathrm{active}}}.
\]
It measures distance below the chosen threshold on the active calibration score
scale. It is not a calibrated attack probability, causal attribution, or
attack-family classifier. These explanations are useful precisely because
their limited meaning is explicit.
